\question{نمونه‌ی تئوری (فرمول، اثبات، فهرست‌ها)}

این فایل فقط یک الگوی ساده است تا ببینید چگونه می‌توانید «تئوری»‌ها را بنویسید.
برای نوشتن متن انگلیسی وسط متن فارسی از \lr{...} استفاده کنید؛ مثلاً \lr{graph search}.

\subsection{فرمول‌ها}
می‌توانید فرمول‌ها را درون متن مثل \(a^2+b^2=c^2\) یا به‌صورت جداگانه بنویسید:
\[
	\sum_{i=1}^{n} i = \frac{n(n+1)}{2}
\]

\subsection{کمی میان‌برهای ریاضی}
مجموعه‌ها و چند نماد پرکاربرد آماده‌اند:
\[
	\IN \subseteq \IZ \subseteq \IQ \subseteq \IR \subseteq \IC
\]
و برای متمم و اندازه‌ی مجموعه:
\[
	\setcomp{A} \provided x \in A \qquad\qquad \card{A}
\]

\subsection{نمونه‌ی یک تعریف و اثبات}
یک تعریف کوتاه:
\[
	f(n) = \polylog(n)
\]

\begin{proof}
	این فقط یک «نمونه‌ی ظاهری» از محیط اثبات است. شما می‌توانید مراحل را با یک فهرست هم بنویسید:
	\begin{ابجد}
		\item بیان فرض‌ها
		\item تبدیل‌ها و نامساوی‌ها
		\item نتیجه‌گیری
	\end{ابجد}
\end{proof}

\subsection{یک نامساویِ تمیز}
در این قالب، \(\leq\) و \(\geq\) به شکل مایل (\(\leqslant,\geqslant\)) نمایش داده می‌شوند:
\[
	\floor{x} \le x \le \ceil{x}
\]

\subsection{استفاده از \lr{\texttt{\textbackslash question}}}
می‌توانید شماره‌ی سوال را به صورت دستی هم تعیین کنید:

\question[2]{یک سوال دوم (نمونه)}
متن پاسخ اینجا قرار می‌گیرد.

